
\documentclass[10pt,a4paper,fontset=none]{ctexart}
\usepackage[margin=1.65cm]{geometry}
\usepackage{xcolor}
\usepackage{listings}
\usepackage{booktabs,longtable}
\usepackage{enumitem}
\usepackage{hyperref}
\usepackage{fancyhdr}
\usepackage{titlesec}
\usepackage{tcolorbox}
\tcbuselibrary{breakable}
\usepackage{fontspec}
\setCJKmainfont{Songti SC}
\setCJKsansfont{Songti SC}
\setCJKmonofont{Songti SC}
\setmonofont{Menlo}[Scale=0.82]
\hypersetup{colorlinks=true, linkcolor=blue!50!black, urlcolor=blue!60!black}
\setlength{\headheight}{14pt}
\pagestyle{fancy}
\fancyhf{}
\lhead{DSA 期末错题总结}
\rhead{按对话代码来源校订版}
\cfoot{\thepage}
\titleformat{\section}{\Large\bfseries\color{blue!45!black}}{\thesection}{0.7em}{}
\titleformat{\subsection}{\large\bfseries}{\thesubsection}{0.6em}{}
\definecolor{codebg}{RGB}{248,248,248}
\definecolor{darkgreen}{RGB}{0,100,0}
\lstdefinestyle{py}{
  language=Python,
  backgroundcolor=\color{codebg},
  basicstyle=\ttfamily\scriptsize,
  keywordstyle=\color{blue!70!black}\bfseries,
  commentstyle=\color{darkgreen},
  stringstyle=\color{red!55!black},
  showstringspaces=false,
  frame=single,
  rulecolor=\color{black!20},
  breaklines=true,
  breakatwhitespace=false,
  tabsize=4,
  columns=fullflexible,
  keepspaces=true
}
\lstdefinestyle{sample}{
  backgroundcolor=\color{codebg},
  basicstyle=\ttfamily\scriptsize,
  frame=single,
  rulecolor=\color{black!15},
  breaklines=true,
  columns=fullflexible,
  keepspaces=true
}
\newtcolorbox{taskbox}{colback=blue!3!white,colframe=blue!35!black,title=题干与样例,breakable}
\newtcolorbox{mistakebox}{colback=red!2!white,colframe=red!40!black,title=错误梳理,breakable}
\newtcolorbox{notebox}{colback=green!2!white,colframe=green!35!black,title=复盘记忆点,breakable}
\newcommand{\codesource}{\par\noindent\textbf{代码来源：}本题代码优先沿用你在对话中贴出的最终/可用版本；如果当前对话记录中没有完整用户最终代码，则采用此前对话中我给出的可用版本。本次校订不重新设计新的算法实现。\par\vspace{0.2em}}
\setlist[itemize]{itemsep=0.16em, topsep=0.25em, leftmargin=1.6em}
\title{\bfseries DSA 期末错题总结：按对话代码来源校订版\\\large 题干样例、错误梳理、对话来源代码与考前模板}
\author{依据你上传的《题库.docx》和项目对话整理}
\date{2026 年 6 月}
\begin{document}
\maketitle
\tableofcontents
\newpage

\section*{使用说明}
\addcontentsline{toc}{section}{使用说明}
这份最终版严格沿用你上传的题库 Word 文档中的原始题号与题名。例如，树题使用 \textbf{22:根据二叉树前中序序列建树}，图题中也保留题库里的重复编号，例如 \textbf{1:单词序列} 和 \textbf{1:拓扑排序}。章节结构仍按复习逻辑分为“树、二叉树、BST 与堆”和“图、BFS、最短路、生成树与拓扑”。每道题不再堆放旧错误代码，而是把你多轮修改中暴露出来的问题合并成简明复盘点。本版另外校订了“最终代码”的来源：不再把本次重新写出的代码当作最终代码；所有代码均按对话来源保留，能确认是你贴出的版本则优先采用，不能确认用户最终版的题目采用此前对话中已经给出的可用代码，并在代码前统一说明。

\section{树、二叉树、BST 与堆}

\subsection{22: 根据二叉树前中序序列建树}
\begin{taskbox}
\textbf{题目编号与名称：}22: 根据二叉树前中序序列建树\par
\textbf{要求：}已知二叉树的前序遍历和中序遍历，结点为互不相同的大写字母，输出后序遍历。多组数据，每组两行。\par
\textbf{输入：}多组数据；每组第一行是前序序列，第二行是中序序列，长度不超过 26。\par
\textbf{输出：}对每组数据输出一行后序遍历序列。\par
\textbf{样例输入：}
\begin{lstlisting}[style=sample]
DURPA
RUDPA
XTCNB
CTBNX
\end{lstlisting}
\textbf{样例输出：}
\begin{lstlisting}[style=sample]
RUAPD
CBNTX
\end{lstlisting}
\end{taskbox}
\begin{mistakebox}
\begin{itemize}
  \item 不要真的先建完整结点对象也可以；本题只要求输出后序，递归字符串切片最直接。
  \item 前序第一个字符一定是根；在中序中找到根，左边是左子树，右边是右子树。
  \item 切分前序时要用左子树长度，不能按字符位置随便截。
  \item 这是多组数据，不能只读两行；要循环读到 EOF。
  \item 递归出口是空串返回空串，单结点也会自然返回根。
\end{itemize}
\end{mistakebox}
\codesource
\begin{lstlisting}[style=py,caption={最终可用代码（对话来源）}]
class Treenode:
    def __init__(self, val = 0, left = None, right = None):
        self.val = val
        self.left = left
        self.right = right

def Buildtree(preorder, inorder):
    if not preorder:
        return None

    root_val = preorder[0]
    root = Treenode(root_val)

    x = inorder.index(root_val)

    # 中序遍历：左子树 根 右子树
    ino_left = inorder[:x]
    ino_right = inorder[x + 1:]

    # 左子树结点个数
    l = len(ino_left)

    # 前序遍历：根 左子树 右子树
    pre_left = preorder[1:1 + l]
    pre_right = preorder[1 + l:]

    root.left = Buildtree(pre_left, ino_left)
    root.right = Buildtree(pre_right, ino_right)

    return root

def postorder(root,ans):
    if root:
        postorder(root.left, ans)
        postorder(root.right, ans)
        ans.append(root.val)

while True:
    try:
        line1 = input()
        line2 = input()
        root = Buildtree(line1, line2)
        ans = []
        postorder(root, ans)
        print(''.join(ans))
    except EOFError:
        break
\end{lstlisting}
\begin{notebox}
\begin{itemize}
  \item 前序 + 中序：根在前序开头，靠中序划分左右子树。
  \item 输出后序时顺序是 left + right + root。
  \item 多组到 EOF 的题，用 sys.stdin 读完整输入更稳。
\end{itemize}
\end{notebox}


\subsection{23: 括号嵌套二叉树}
\begin{taskbox}
\textbf{题目编号与名称：}23: 括号嵌套二叉树\par
\textbf{要求：}给出括号嵌套表示的二叉树，输出前序遍历和中序遍历。星号 * 表示空树；形如 A(B(*,C),D(E))。\par
\textbf{输入：}第一行 n；接下来 n 行，每行一棵树的括号嵌套表示。\par
\textbf{输出：}对每棵树输出两行：前序遍历和中序遍历。\par
\textbf{样例输入：}
\begin{lstlisting}[style=sample]
2
A
A(B(*,C),D(E))
\end{lstlisting}
\textbf{样例输出：}
\begin{lstlisting}[style=sample]
A
A
ABCDE
BCAED
\end{lstlisting}
\end{taskbox}
\begin{mistakebox}
\begin{itemize}
  \item 这题最容易错在解析括号：不能简单按逗号 split，因为逗号可能出现在更深层括号里。
  \item 根一定是当前子串第一个字符；如果后面没有括号，就是叶子。
  \item 左子树为空、右子树非空时写作 A(*,R)，不能把 * 当普通结点输出。
  \item 题面文字提到后序，但输出要求只有前序和中序，要以输出要求为准。
  \item 多棵树逐行处理，不要把全部输入拼成一个大字符串。
\end{itemize}
\end{mistakebox}
\codesource
\begin{lstlisting}[style=py,caption={最终可用代码（对话来源）}]
class Treenode:
    def __init__(self, val = 0, left = None, right = None):
        self.val = val
        self.left = left
        self.right = right
def buildtree(s,i):
    if s[i] == "*":
        return None, i+1
    root = Treenode(s[i])
    i += 1
    if i < len(s) and s[i] == "(":
        i+=1
        root.left, i = buildtree(s, i)
        if i < len(s) and s[i] == ",":
            i += 1
            root.right, i = buildtree(s,i)

        if i < len(s) and s[i] == ")":
            i+=1

    return root,i

def preorder(root, ans):
    if root:
        ans.append(root.val)
        preorder(root.left, ans)
        preorder(root.right, ans)

def inorder(root, ans):
    if root:
        inorder(root.left, ans)
        ans.append(root.val)
        inorder(root.right, ans)

n = int(input())
for i in range(n):
    s = input()
    root, j= buildtree(s,0)
    ans1 = []
    ans2 = []
    preorder(root, ans1)
    inorder(root, ans2)
    print("".join(ans1))
    print("".join(ans2))
\end{lstlisting}
\begin{notebox}
\begin{itemize}
  \item 括号解析题的关键是“只切最外层逗号”。
  \item 遇到 * 返回 None，遍历时 None 输出空串。
  \item 输出什么遍历要看“输出”部分，不要被描述里多余的话带偏。
\end{itemize}
\end{notebox}


\subsection{24: 二叉树}
\begin{taskbox}
\textbf{题目编号与名称：}24: 二叉树\par
\textbf{要求：}完全二叉树编号为 1,2,3,...,n，求编号 m 所在子树中实际存在的结点个数。\par
\textbf{输入：}多组数据，每行两个整数 m,n；0 0 结束。\par
\textbf{输出：}每组输出一行，表示以 m 为根的子树中编号不超过 n 的结点个数。\par
\textbf{样例输入：}
\begin{lstlisting}[style=sample]
3 12
0 0
\end{lstlisting}
\textbf{样例输出：}
\begin{lstlisting}[style=sample]
4
\end{lstlisting}
\end{taskbox}
\begin{mistakebox}
\begin{itemize}
  \item 不要建树：n 最大到 1e9，建结点一定不合适。
  \item 完全二叉树中，m 的下一层不是一个结点，而是连续区间 [2m,2m+1]。
  \item 每一层都要与 n 截断，右端超过 n 时只能统计到 n。
  \item 输入是多组，以 0 0 结束；你之前容易只处理一组。
  \item 题面输入顺序是 m,n，不要写反成 n,m。
\end{itemize}
\end{mistakebox}
\codesource
\begin{lstlisting}[style=py,caption={最终可用代码（对话来源）}]
while True:
    n,m = map(int, input().split())
    if n==0 and m==0:
        break
    ans =0
    right = n
    left = n
    while left <= m:
        ans += min(m, right) - left + 1
        right = right*2 + 1
        left = left*2
    print(ans)
\end{lstlisting}
\begin{notebox}
\begin{itemize}
  \item 堆式编号题优先想“每层区间扩张”。
  \item while 条件看 left 是否仍然存在：left <= n。
  \item 贡献公式：min(right,n)-left+1。
\end{itemize}
\end{notebox}


\subsection{25: 二叉搜索树的层次遍历}
\begin{taskbox}
\textbf{题目编号与名称：}25: 二叉搜索树的层次遍历\par
\textbf{要求：}输入一行无序整数，重复数字只插入一次，建立二叉搜索树后输出层次遍历。\par
\textbf{输入：}一行若干整数，空格分隔。\par
\textbf{输出：}一行输出 BST 的层次遍历结果，数字之间一个空格，末尾无多余空格。\par
\textbf{样例输入：}
\begin{lstlisting}[style=sample]
51 45 59 86 45 4 15 76 60 20 61 77 62 30 2 37 13 82 19 74 2 79 79 97 33 90 11 7 29 14 50 1 96 59 91 39 34 6 72 7
\end{lstlisting}
\textbf{样例输出：}
\begin{lstlisting}[style=sample]
51 45 59 4 50 86 2 15 76 97 1 13 20 60 77 90 11 14 19 30 61 82 96 7 29 37 62 79 91 6 33 39 74 34 72
\end{lstlisting}
\end{taskbox}
\begin{mistakebox}
\begin{itemize}
  \item 重复值只计入一次，不能简单把相等元素插到左边或右边。
  \item 层次遍历要用队列；用 list.pop(0) 可以过小数据，但 deque 更规范。
  \item 输出末尾不能多空格，最好先收集成字符串列表再 join。
  \item 插入函数中 key == node.key 时直接返回，不能继续递归。
  \item 如果输入为空要避免访问 root，不过本题通常保证有数字。
\end{itemize}
\end{mistakebox}
\codesource
\begin{lstlisting}[style=py,caption={最终可用代码（对话来源）}]
class TreeNode:
    def __init__(self, key):
        self.key = key
        self.left = None
        self.right = None

class BinarySearchTree:
    def __init__(self):
        self.root = None
    def put(self,key):
        if self.root:
            self._put(key, self.root)
        else:
            self.root = TreeNode(key)
    def _put(self, key, currentNode):
        if key < currentNode.key:
            if currentNode.left:
                self._put(key,currentNode.left)
            else:
                currentNode.left = TreeNode(key)
        elif key > currentNode.key:
            if currentNode.right:
                self._put(key, currentNode.right)
            else:
                currentNode.right = TreeNode(key)
def bfs(root,ans):
    if not root:
        return
    else:
        queue = [root]
        while queue:
            node = queue.pop(0)
            ans.append(str(node.key))
            if node.left:
                queue.append(node.left)
            if node.right:
                queue.append(node.right)
    return ans

lis = list(map(int,input().split()))
hp = BinarySearchTree()
for n in lis:
    hp.put(n)
ans = []
bfs(hp.root, ans)
print(" ".join(ans))

\end{lstlisting}
\begin{notebox}
\begin{itemize}
  \item BST 插入：小走左，大走右，相等忽略。
  \item 层次遍历模板：deque + popleft。
  \item 输出列表用 join，避免末尾空格。
\end{itemize}
\end{notebox}


\subsection{26: 实现堆结构}
\begin{taskbox}
\textbf{题目编号与名称：}26: 实现堆结构\par
\textbf{要求：}维护一个初始为空的数组，支持插入元素和输出删除最小元素，要求用最小堆高效实现。\par
\textbf{输入：}第一行 t；每组第一行 n；之后 n 行操作。type=1 后接 u 表示插入，type=2 表示输出并删除最小值。\par
\textbf{输出：}每次 type=2 输出被删除的最小数字。\par
\textbf{样例输入：}
\begin{lstlisting}[style=sample]
2
5
1 1
1 2
1 3
2
2
4
1 5
1 1
1 7
2
\end{lstlisting}
\textbf{样例输出：}
\begin{lstlisting}[style=sample]
1
2
1
\end{lstlisting}
\end{taskbox}
\begin{mistakebox}
\begin{itemize}
  \item Python 模块名是 heapq，不是 headq；弹出函数是 heappop，不是 heapop。
  \item 每组测试数据都要重新建一个空堆，不要把上一组残留带到下一组。
  \item 操作行长度不固定：type=1 有两个数，type=2 只有一个数。
  \item 复杂度要求 O(n log n)，不能每次删除前 sort。
  \item 输出多行建议收集到 ans 最后统一打印，速度更稳。
\end{itemize}
\end{mistakebox}
\codesource
\begin{lstlisting}[style=py,caption={最终可用代码（对话来源）}]
class BinHeap:
    def __init__(self):
        self.heaplist = [0]
        self.size = 0
    def percUp(self, i):
        while i//2 > 0:
            if self.heaplist[i] < self.heaplist[i//2]:
                self.heaplist[i], self.heaplist[i//2] = self.heaplist[i//2], self.heaplist[i]
            i = i//2
    def insert(self, k):
        self.heaplist.append(k)
        self.size += 1
        self.percUp(self.size)
    def FindMin(self,i):
        if i*2+1>self.size:
            return i*2
        else:
            if self.heaplist[i*2] < self.heaplist[i*2+1]:
                return i*2
            else:
                return i*2+1
    def percDown(self,i):
        while i*2 <= self.size:
            mc = self.FindMin(i)
            if self.heaplist[i] > self.heaplist[mc]:
                self.heaplist[i], self.heaplist[mc] = self.heaplist[mc], self.heaplist[i]
            i = mc
    def delMin(self):
        val = self.heaplist[1]
        self.heaplist[1] = self.heaplist[self.size] # 注意这里size和实际队列长度的区别
        self.size -= 1
        self.heaplist.pop()
        self.percDown(1) # 这里是要往下沉的所以是1
        return val
    
t = int(input())

while t:
    hp = BinHeap()
    n = int(input())

    while n:
        op = list(map(int, input().split()))

        if op[0] == 1:
            hp.insert(op[1])
        else:
            print(hp.delMin())

        n -= 1

    t -= 1
\end{lstlisting}
\begin{notebox}
\begin{itemize}
  \item heapq 默认就是最小堆。
  \item 插入 heappush，删除最小 heappop。
  \item 多组测试时，每组 heap=[]。
\end{itemize}
\end{notebox}


\subsection{27: 树的转换}
\begin{taskbox}
\textbf{题目编号与名称：}27: 树的转换\par
\textbf{要求：}用 d/u 字符串表示一般树的 DFS 过程，将其按左儿子右兄弟法转换为二叉树，输出转换前和转换后的高度。\par
\textbf{输入：}一行由 u 和 d 组成的字符串。\par
\textbf{输出：}按格式 h1 => h2 输出。\par
\textbf{样例输入：}
\begin{lstlisting}[style=sample]
dudduduudu
\end{lstlisting}
\textbf{样例输出：}
\begin{lstlisting}[style=sample]
2 => 4
\end{lstlisting}
\end{taskbox}
\begin{mistakebox}
\begin{itemize}
  \item 普通树高度只看向下 d 的最大深度。
  \item 转换后二叉树高度中，向第一个孩子走和向右兄弟走都会增加边数。
  \item 你之前容易把“兄弟边”不算进转换后二叉树高度，导致样例对不上。
  \item 这题可以不真的建树，直接 DFS 模拟时同步维护两个高度更简单。
  \item 输出格式固定，=> 两边有空格。
\end{itemize}
\end{mistakebox}
\codesource
\begin{lstlisting}[style=py,caption={最终可用代码（对话来源）}]
class TreeNode:
    def __init__(self, key):
        self.key = key
        self.children = []
        self.parent = None   # 普通树：可以有多个孩子


class BinaryNode:
    def __init__(self, key):
        self.key = key
        self.left = None     # 左儿子
        self.right = None # 右兄弟


def convert(root):
    if root is None:
        return None

    # 创建对应的二叉树结点
    b_root = BinaryNode(root.key)

    # 如果没有孩子，直接返回
    if len(root.children) == 0:
        return b_root

    # 第一个孩子变成左儿子
    b_root.left = convert(root.children[0])

    # 其余孩子依次变成右兄弟
    current = b_root.left

    for child in root.children[1:]:
        current.right = convert(child)
        current = current.right

    return b_root

def height(root):
    if root is None:
        return -1
    else:
        height_left = height(root.left)
        height_right = height(root.right)
        return max(height_left, height_right) + 1

def length(root):
    if root is None:
        return -1
    if len(root.children) == 0:
        return 0
    max_height = 0
    for child in root.children:
        max_height = max(max_height, length(child))
    return max_height + 1

n = input().strip()
count = 0
root = TreeNode(0)
current = root
for i in n:
    if i == 'd':
        count += 1
        nd = TreeNode(count)
        current.children.append(nd)
        nd.parent = current
        current = nd
    elif i == 'u':
        current = current.parent

b_root = convert(root)
print(f"{length(root)} => {height(b_root)}")
\end{lstlisting}
\begin{notebox}
\begin{itemize}
  \item 原树高度：最大向下层数。
  \item 左儿子右兄弟二叉树高度：child 和 sibling 边都算。
  \item 这题不用输出树，只要维护高度。
\end{itemize}
\end{notebox}


\subsection{28: Huffman编码树}
\begin{taskbox}
\textbf{题目编号与名称：}28: Huffman编码树\par
\textbf{要求：}给定 n 个外部结点权值，构造 Huffman 树，求最小带权外部路径长度。\par
\textbf{输入：}第一行 t；每组第一行 n，第二行 n 个权值。\par
\textbf{输出：}每组输出一行最小外部路径长度总和。\par
\textbf{样例输入：}
\begin{lstlisting}[style=sample]
2
3
1 2 3
4
1 1 3 5
\end{lstlisting}
\textbf{样例输出：}
\begin{lstlisting}[style=sample]
9
17
\end{lstlisting}
\end{taskbox}
\begin{mistakebox}
\begin{itemize}
  \item Huffman 的核心不是排序一次，而是每次取当前最小的两个合并。
  \item 合并产生的新权值要重新放回集合中继续参与比较。
  \item 答案累加每次合并的和，而不是最后根结点权值。
  \item 即使 n 小，也建议用堆；手动列表删除容易写错索引。
  \item 多组数据时每组重新 heapify。
\end{itemize}
\end{mistakebox}
\codesource
\begin{lstlisting}[style=py,caption={最终可用代码（对话来源）}]
import heapq

def optimal_binary_tree(weights):
    # 1. 开始 n 个结点位于集合 S
    heapq.heapify(weights)

    total_wpl = 0

    # 2. 每次取出两个权值最小的结点
    while len(weights) > 1:
        n1 = heapq.heappop(weights)
        n2 = heapq.heappop(weights)

        # 构造新结点 r
        r = n1 + n2

        # 这一次合并产生的代价
        total_wpl += r

        # 将 r 加回集合 S
        heapq.heappush(weights, r)

    # 3. total_wpl 就是最小 WPL
    return total_wpl

t = int(input())
for _ in range(t):
    n = int(input())
    weights = list(map(int, input().split()))  # 这一行读入 n 个权值
    print(optimal_binary_tree(weights))
\end{lstlisting}
\begin{notebox}
\begin{itemize}
  \item Huffman 模板：反复取最小两个，合并，答案加合并值。
  \item heapify 是把普通列表原地变成堆。
  \item 答案不是编码字符串，而是最小 WPL。
\end{itemize}
\end{notebox}


\subsection{49: 二叉树的深度}
\begin{taskbox}
\textbf{题目编号与名称：}49: 二叉树的深度\par
\textbf{要求：}给定编号 1 到 n 的二叉树，根结点为 1，求深度。深度按根到叶最长路径上的结点数计算。\par
\textbf{输入：}第一行 n；接下来 n 行，第 i 行给出结点 i 的左、右孩子编号，-1 表示空。\par
\textbf{输出：}输出一个整数，即二叉树深度。\par
\textbf{样例输入：}
\begin{lstlisting}[style=sample]
3
2 3
-1 -1
-1 -1
\end{lstlisting}
\textbf{样例输出：}
\begin{lstlisting}[style=sample]
2
\end{lstlisting}
\end{taskbox}
\begin{mistakebox}
\begin{itemize}
  \item 这题根明确是 1，不需要再找根。
  \item 输入行号就是结点编号：第 i 行描述结点 i，不能按遍历顺序读。
  \item 深度按结点数，不是边数。空树为 0，叶子为 1。
  \item n 很小，但写成通用递归更安全。
  \item 左右孩子编号为 -1 时不要访问 children[-1]。
\end{itemize}
\end{mistakebox}
\codesource
\begin{lstlisting}[style=py,caption={最终可用代码（对话来源）}]
class BinaryNode:
    def __init__(self, key):
        self.key = key
        self.left = None     # 左儿子
        self.right = None # 右兄弟
def height(root):
    if root is None:
        return 0
    else:
        height_left = height(root.left)
        height_right = height(root.right)
        return max(height_left, height_right) + 1

n = int(input())
nodes = [None] * (n + 1)
for i in range(1, n + 1):
    nodes[i] = BinaryNode(i)
for i in range(1, n + 1):
    a, b = map(int, input().split())
    if a != -1:
        nodes[i].left = nodes[a]
    if b != -1:
        nodes[i].right = nodes[b]
root = nodes[1]

print(height(root))
\end{lstlisting}
\begin{notebox}
\begin{itemize}
  \item 如果题目写“根结点为 1”，直接 depth(1)。
  \item 数组 left/right 比建对象更短。
  \item 递归出口：编号 -1 深度 0。
\end{itemize}
\end{notebox}


\subsection{50: 扩展二叉树}
\begin{taskbox}
\textbf{题目编号与名称：}50: 扩展二叉树\par
\textbf{要求：}给出扩展二叉树的先序序列，其中 . 表示空结点，输出原二叉树的中序和后序。\par
\textbf{输入：}一行扩展二叉树先序序列，由大写字母和 . 组成。\par
\textbf{输出：}第一行输出中序序列，第二行输出后序序列。\par
\textbf{样例输入：}
\begin{lstlisting}[style=sample]
ABD..EF..G..C..
\end{lstlisting}
\textbf{样例输出：}
\begin{lstlisting}[style=sample]
DBFEGAC
DFGEBCA
\end{lstlisting}
\end{taskbox}
\begin{mistakebox}
\begin{itemize}
  \item 扩展先序可以唯一建树，关键是读到 . 时返回空。
  \item 需要一个全局下标或迭代器，不能每次从字符串开头重新读。
  \item 输出时不要输出 .，它只表示空结点。
  \item 中序是 left-root-right，后序是 left-right-root，顺序不要混。
  \item 递归构建时必须先建左子树，再建右子树，因为输入是先序。
\end{itemize}
\end{mistakebox}
\codesource
\begin{lstlisting}[style=py,caption={最终可用代码（对话来源）}]
class Treenode:
    def __init__(self, val = 0, left = None, right = None):
        self.val = val
        self.left = left
        self.right = right

preorder = input().strip()

def build_tree(seq):
    """根据扩展二叉树的先序序列递归建树"""
    idx = 0  # 用于记录当前处理的位置
    def helper():
        nonlocal idx
        if idx >= len(seq):
            return None

        val = seq[idx]
        idx += 1

        # 遇到 '.' 表示空节点
        if val == '.':
            return None

        # 创建当前节点
        node = Treenode(val)

        # 递归构建左子树
        node.left = helper()
        # 递归构建右子树
        node.right = helper()

        return node

    return helper()

def inorder(root, ans):
    if root:
        inorder(root.left,ans)
        ans.append(root.val)
        inorder(root.right,ans)
    return ans

def postorder(root, ans):
    if root:
        postorder(root.left,ans)
        postorder(root.right,ans)
        ans.append(root.val)
    return ans


# 构建二叉树
root = build_tree(preorder)
ans1 = []
ans2 = []
ans1 = inorder(root, ans1)
ans2 = postorder(root, ans2)
print(''.join(ans1))
print(''.join(ans2))

\end{lstlisting}
\begin{notebox}
\begin{itemize}
  \item 扩展二叉树题：看到 . 就返回 None。
  \item 先序建树顺序：root, left, right。
  \item 遍历输出时忽略空结点。
\end{itemize}
\end{notebox}


\subsection{53: 根据二叉树中后序序列建树}
\begin{taskbox}
\textbf{题目编号与名称：}53: 根据二叉树中后序序列建树\par
\textbf{要求：}已知中序遍历和后序遍历，结点互不相同，输出前序遍历。\par
\textbf{输入：}两行：第一行中序序列，第二行后序序列。\par
\textbf{输出：}输出一行前序遍历。\par
\textbf{样例输入：}
\begin{lstlisting}[style=sample]
DUPR
DPRU
\end{lstlisting}
\textbf{样例输出：}
\begin{lstlisting}[style=sample]
UDRP
\end{lstlisting}
\end{taskbox}
\begin{mistakebox}
\begin{itemize}
  \item 后序最后一个字符是根，不是第一个。
  \item 在中序中找到根后，左边是左子树，右边是右子树。
  \item 切分后序时仍然要靠左子树长度，右子树在根之前。
  \item 输出前序是 root + left + right。
  \item 长度不超过 26，字符串递归足够。
\end{itemize}
\end{mistakebox}
\codesource
\begin{lstlisting}[style=py,caption={最终可用代码（对话来源）}]
class Treenode:
    def __init__(self, val = 0, left = None, right = None):
        self.val = val
        self.left = left
        self.right = right

def Buildtree(postorder, inorder):
    if not postorder: # 这个判断是非常有意义的
        return 
    root_val = postorder[-1] # 赋值
    root = Treenode(root_val) # 赋一个结点

    x = inorder.index(root_val)

    ino_left = inorder[:x]
    ino_right = inorder[x+1:]

    l = len(ino_left)

    post_left = postorder[0:l]
    post_right = postorder[l:-1]

    root.left = Buildtree(post_left, ino_left)
    root.right = Buildtree(post_right, ino_right)

    return root

def preorder(root, ans):
    if root:
        ans.append(root.val)
        preorder(root.left,ans)
        preorder(root.right,ans)
    return ans

io = input()
po = input()
ans = []
root = Buildtree(po,io)

ans = preorder(root,ans)

print("".join(ans))
\end{lstlisting}
\begin{notebox}
\begin{itemize}
  \item 中序 + 后序：根在后序末尾。
  \item 靠中序划分左右，靠左子树长度切后序。
  \item 输出前序：root + left + right。
\end{itemize}
\end{notebox}


\subsection{54: 验证二叉搜索树}
\begin{taskbox}
\textbf{题目编号与名称：}54: 验证二叉搜索树\par
\textbf{要求：}给定带 \# 占位的层次遍历序列，判断它是否能构成合法 BST。\par
\textbf{输入：}一行若干字符串，\# 表示空结点，非空结点为唯一整数。\par
\textbf{输出：}合法输出 YES，否则输出 NO。\par
\textbf{样例输入：}
\begin{lstlisting}[style=sample]
10 5 15 2 7 12 20
\end{lstlisting}
\textbf{样例输出：}
\begin{lstlisting}[style=sample]
YES
\end{lstlisting}
\end{taskbox}
\begin{mistakebox}
\begin{itemize}
  \item 不能只判断每个结点与左右孩子的大小关系；必须检查整棵子树的取值范围。
  \item 例如右子树中的所有点都必须大于根，不只是右孩子大于根。
  \item \# 是占位，不要转成整数。
  \item 层次数组中下标 i 的孩子是 2i+1 和 2i+2。
  \item 边界必须严格：左 < root < 右，不能允许等号。
\end{itemize}
\end{mistakebox}
\codesource
\begin{lstlisting}[style=py,caption={最终可用代码（对话来源）}]
from collections import deque

# 定义节点
class Node:
    def __init__(self, val):
        self.val = val
        self.left = None
        self.right = None

# 构建二叉树
def build_tree_level(level_list):
    if not level_list:
        return None
    root_val = level_list[0]
    root = Node(int(root_val))
    queue = deque([root])
    i = 1
    while queue and i < len(level_list):
        node = queue.popleft()
        # 左孩子
        if i < len(level_list):
            if level_list[i] != '#':
                node.left = Node(int(level_list[i]))
                queue.append(node.left)
            i += 1
        # 右孩子
        if i < len(level_list):
            if level_list[i] != '#':
                node.right = Node(int(level_list[i]))
                queue.append(node.right)
            i += 1
    return root

# 验证 BST
def is_bst(node, low=-float('inf'), high=float('inf')):
    if node is None:
        return True
    val = node.val
    if val <= low or val >= high:
        return False
    return is_bst(node.left, low, val) and is_bst(node.right, val, high)

# 主程序
level_input = input().split()
root = build_tree_level(level_input)

if is_bst(root):
    print("YES")
else:
    print("NO")
\end{lstlisting}
\begin{notebox}
\begin{itemize}
  \item BST 验证的本质是“范围传递”，不是只看父子。
  \item 左子树上界变成根值，右子树下界变成根值。
  \item 层次遍历带 \# 时可以直接用数组下标递归。
\end{itemize}
\end{notebox}

\section{图、BFS、最短路、生成树与拓扑}

\subsection{1: 单词序列}
\begin{taskbox}
\textbf{题目编号与名称：}1: 单词序列\par
\textbf{要求：}给出开始单词、结束单词和词典，每次只能改一个字母，中间单词必须在词典中，求最短转换序列长度。\par
\textbf{输入：}第一行开始单词和结束单词；第二行词典单词，空格分隔。\par
\textbf{输出：}输出最短转换序列长度；无法转换输出 0。\par
\textbf{样例输入：}
\begin{lstlisting}[style=sample]
hit cog
hot dot dog lot log
\end{lstlisting}
\textbf{样例输出：}
\begin{lstlisting}[style=sample]
5
\end{lstlisting}
\end{taskbox}
\begin{mistakebox}
\begin{itemize}
  \item 这是无权最短路，应该用 BFS，不是 DFS。
  \item 开始单词和结束单词可以不在词典中，所以要把它们加入候选点集合。
  \item 只差一个字母才能连边，不能用字符串相似度或排序判断。
  \item 距离按序列长度算，hit->...->cog 包含起点和终点，所以起点距离为 1。
  \item visited 要在入队时标记，避免同一层反复入队。
\end{itemize}
\end{mistakebox}
\codesource
\begin{lstlisting}[style=py,caption={最终可用代码（对话来源）}]
from collections import deque

class Vertex:
    def __init__(self, key):
        self.id = key
        self.connectedTo = {}
        self.color = 'white' # 顶点访问状态（white=未访问，gray=正在访问，black=已完成）
        self.distance = float('inf') # 当前顶点到起始点的距离，初始设为无穷大
        self.pred = None # 处理这个顶点的前驱，用于记录路径
 
    def addNeighbor(self, nbr, weight=0):
        self.connectedTo[nbr] = weight # 是字典，nbr是顶点对象的key

    def __str__(self):
        return str(self.id) + ' connectedTo: ' + str([x.id for x in self.connectedTo])

    def getConnections(self):
        return self.connectedTo.keys()

    def getId(self):
        return self.id

    def getWeight(self, nbr):
        return self.connectedTo[nbr]

    def getPred(self):
        return self.pred

    def setPred(self, pred):
        self.pred = pred

    def getColor(self):
        return self.color

    def setColor(self, color):
        self.color = color

    def getDistance(self):
        return self.distance

    def setDistance(self, distance):
        self.distance = distance


class Graph:

    def __init__(self):
        self.vertList = {}
        self.numVertices = 0

    def addVertex(self, key):
        self.numVertices = self.numVertices + 1
        newVertex = Vertex(key)
        self.vertList[key] = newVertex
        return newVertex

    def getVertex(self, n): # 通过key查找顶点
        if n in self.vertList:
            return self.vertList[n]
        else:
            return None

    def __contains__(self, n):
        return n in self.vertList
    
    def addEdge(self, f, t, cost=0): # f起点，t终点
        if f not in self.vertList:
            nv = self.addVertex(f)

        if t not in self.vertList:
            nv = self.addVertex(t)
        self.vertList[f].addNeighbor(self.vertList[t], cost)

    def getVertices(self):
        return self.vertList.keys()

    def __iter__(self):
        return iter(self.vertList.values())


def buildGraph(wordlist):
    d = {}
    g = Graph()

    # 第一步：建立 bucket
    # bucket 的作用：把“只差一个字母”的单词放到同一组里
    for word in wordlist:

        # 依次把单词的每一个位置替换成 "_"
        # 例如 fool:
        # _ool, f_ol, fo_l, foo_
        for i in range(len(word)):
            bucket = word[:i] + '_' + word[i + 1:]

            # 如果这个 bucket 已经存在，就把当前单词加入这个 bucket
            if bucket in d:
                d[bucket].append(word)

            # 如果这个 bucket 不存在，就创建一个新的列表
            else:
                d[bucket] = [word]

    # 第二步：根据 bucket 建图
    # 同一个 bucket 中的单词，说明它们只差一个字母
    # 因此这些单词之间应该有边
    for bucket in d.keys():
        for word1 in d[bucket]:
            for word2 in d[bucket]:

                # 避免自己和自己连边
                if word1 != word2:
                    g.addEdge(word1, word2)

    # 返回构建好的图
    return g

def bfs(g, start):
    # 将起始顶点到自己的距离设置为 0
    start.setDistance(0)

    # 起始顶点没有前驱节点
    start.setPred(None)

    start.setColor('gray')

    # 创建一个队列，用于 BFS 的“先进先出”遍历
    vertQueue = deque()

    # 先把起始顶点放入队列
    vertQueue.append(start)

    # 只要队列不为空，就继续搜索
    while len(vertQueue) > 0:

        # 取出队首顶点，作为当前正在访问的顶点
        currentVert = vertQueue.popleft()

        # 遍历当前顶点的所有邻接点
        for nbr in currentVert.getConnections():

            # 如果邻接点还是 white，说明它还没有被访问过
            if nbr.getColor() == 'white':

                # 将邻接点标记为 gray，表示已经发现，但还没有完全处理完
                nbr.setColor('gray')

                # 邻接点距离 = 当前顶点距离 + 1
                # 因为 BFS 每走一条边，距离增加 1
                nbr.setDistance(currentVert.getDistance() + 1)

                # 记录邻接点的前驱节点
                # 说明 nbr 是从 currentVert 访问到的
                nbr.setPred(currentVert)

                # 把这个邻接点加入队列，等待之后继续访问它的邻居
                vertQueue.append(nbr)

        # 当前顶点的所有邻居都检查完后，将其标记为 black
        # black 表示这个顶点已经彻底处理完成
        currentVert.setColor('black')

start, end = input().split()
wordlist = input().split()

wordlist.append(start)
wordlist.append(end)

wordlist = list(set(wordlist))

g = buildGraph(wordlist)

s = g.getVertex(start)
e = g.getVertex(end)

if s is None or e is None:
    print(0)
else:
    bfs(g, s)

    if e.getDistance() == float('inf'):
        print(0)
    else:
        print(e.getDistance() + 1)
\end{lstlisting}
\begin{notebox}
\begin{itemize}
  \item “最短转换次数/序列长度”通常用 BFS。
  \item 起点到起点的序列长度是 1。
  \item 单词数不超过 30，O(n\textasciicircum{}2 * L) 建图或直接枚举都能过。
\end{itemize}
\end{notebox}


\subsection{2: 仙岛求药}
\begin{taskbox}
\textbf{题目编号与名称：}2: 仙岛求药\par
\textbf{要求：}在 M*N 迷宫中从 @ 走到 *，\# 不可走，. 可走，求最少经过方格数；包括起点方格。多组数据，0 0 结束。\par
\textbf{输入：}每组先读 M,N；接下来 M 行地图。\par
\textbf{输出：}每组输出最少方格数；不可达输出 -1。\par
\textbf{样例输入：}
\begin{lstlisting}[style=sample]
8 8
.@##...#
#....#.#
#.#.##..
..#.###.
#.#...#.
..###.#.
...#.*..
.#...###
0 0
\end{lstlisting}
\textbf{样例输出：}
\begin{lstlisting}[style=sample]
10
\end{lstlisting}
\end{taskbox}
\begin{mistakebox}
\begin{itemize}
  \item visited 的维度必须是 M 行 N 列，不是 N*N。
  \item 题目要求计数包括初始位置，所以起点步数/距离应从 1 开始。
  \item 多组数据直到 0 0，不能只读一张地图。
  \item 找到 @ 后要记录起点；地图中 * 是终点。
  \item 越界判断用 nx<0 or nx>=M or ny<0 or ny>=N，行列不能写反。
\end{itemize}
\end{mistakebox}
\codesource
\begin{lstlisting}[style=py,caption={最终可用代码（对话来源）}]
import collections

def bfs(maze, M, N, sx,sy):

    # 创建队列，相当于 BFS 中的 Open 表
    q = collections.deque()

    # visited[i] 表示位置 i 是否已经访问过
    # 用于判重，避免反复访问同一个位置
    visited = [[False] * N for i in range(M)]

    directions = [(-1, 0), (1, 0), (0, -1), (0, 1)]

    q.append((sx,sy,0))

    visited[sx][sy] = True

    # 只要队列不为空，就继续 BFS
    while len(q) > 0:
        # 取出队首元素
        x,y,dis = q.popleft()

        # 如果当前位置已经是目标位置 K，输出步数并结束
        if maze[x][y] == '*':
            return dis
        
        for dx, dy in directions:
            nx = x + dx
            ny = y + dy

            if nx < 0 or nx >= M or ny < 0 or ny >= N:
                continue

            if visited[nx][ny] == True:
                continue

            if maze[nx][ny] == '#':
                continue

            visited[nx][ny] = True
            q.append((nx,ny,dis+1))
    return -1

while True:
    M, N = map(int, input().split())

    if M == 0 and N == 0:
        break

    maze = []
    sx, sy = -1, -1

    for i in range(M):
        row = input().strip()
        maze.append(row)

        for j in range(N):
            if row[j] == '@':
                sx, sy = i, j

    print(bfs(maze, M, N, sx, sy))
\end{lstlisting}
\begin{notebox}
\begin{itemize}
  \item 网格 BFS 三件套：队列、visited、四方向。
  \item 先判越界，再判障碍，再判 visited。
  \item 题目说“经过方格数包括起点”时，起点距离设为 1。
\end{itemize}
\end{notebox}


\subsection{3: 变换的迷宫}
\begin{taskbox}
\textbf{题目编号与名称：}3: 变换的迷宫\par
\textbf{要求：}从 S 到 E，每步耗时 1，不能停留；\# 在当前时间为 K 的倍数时可通过，其余时间不可通过，求最短时间。\par
\textbf{输入：}第一行 T；每组 R,C,K；接下来 R 行迷宫。\par
\textbf{输出：}可达输出最短时间，不可达输出 Oop!。\par
\textbf{样例输入：}
\begin{lstlisting}[style=sample]
1
6 6 2
...S..
...#..
.#....
...#..
...#..
..#E#.
\end{lstlisting}
\textbf{样例输出：}
\begin{lstlisting}[style=sample]
7
\end{lstlisting}
\end{taskbox}
\begin{mistakebox}
\begin{itemize}
  \item 状态不能只用 (x,y)，因为同一格在不同 time\%K 下后续可走性不同。
  \item visited 应为 R*C*K 三维，第三维是时间模 K。
  \item 移动到下一格后的时间是 t+1，判断 \# 是否可走要看 next\_time \% K。
  \item 题目说不能停留，所以不能加入原地不动方向。
  \item 输出失败字符串是 Oop!，大小写和感叹号都要准确。
\end{itemize}
\end{mistakebox}
\codesource
\begin{lstlisting}[style=py,caption={最终可用代码（对话来源）}]
import collections

def bfs(maze, M, N, sx,sy,K):

    # 创建队列，相当于 BFS 中的 Open 表
    q = collections.deque()

    # visited[i] 表示位置 i 是否已经访问过
    # 用于判重，避免反复访问同一个位置
    visited = [[[False] * K for i in range(N)]for j in range(M)]

    directions = [(-1, 0), (1, 0), (0, -1), (0, 1)]

    q.append((sx,sy,0))

    visited[sx][sy][0] = True

    # 只要队列不为空，就继续 BFS
    while len(q) > 0:
        # 取出队首元素
        x,y,dis = q.popleft()

        # 如果当前位置已经是目标位置 K，输出步数并结束
        if maze[x][y] == 'E':
            return dis
        
        for dx, dy in directions:
            nx = x + dx
            ny = y + dy

            if nx < 0 or nx >= M or ny < 0 or ny >= N:
                continue

            if visited[nx][ny][(dis+1)%K] == True:
                continue

            if maze[nx][ny] == '#' and (dis+1)%K != 0:
                continue

            visited[nx][ny][(dis+1)%K] = True

            q.append((nx,ny,dis+1))
    return -1

n = int(input())

while n:
    M,N,K = map(int, input().split())
    maze = []
    for i in range(M):
        row = input().strip()
        maze.append(row)

        for j in range(N):
            if row[j] == 'S':
                sx, sy = i, j

    ans= bfs(maze, M, N, sx, sy,K)
    if ans == -1:
        print("Oop!")
    else:
        print(ans)
    n-=1
\end{lstlisting}
\begin{notebox}
\begin{itemize}
  \item 动态迷宫 BFS：位置 + 时间模数 一起作为状态。
  \item 判断石头能不能走，要用“到达那一刻”的时间。
  \item 不能停留，就只有四个移动方向。
\end{itemize}
\end{notebox}


\subsection{1: 拓扑排序}
\begin{taskbox}
\textbf{题目编号与名称：}1: 拓扑排序\par
\textbf{要求：}给出有向图，输出拓扑排序序列；在同等条件下编号小的顶点在前。\par
\textbf{输入：}第一行 v,a；接下来 a 行每行一条弧的两个顶点编号。\par
\textbf{输出：}输出若干个用空格隔开的顶点，格式为 v1 v3 ...。\par
\textbf{样例输入：}
\begin{lstlisting}[style=sample]
6 8
1 2
1 3
1 4
3 2
3 5
4 5
6 4
6 5
\end{lstlisting}
\textbf{样例输出：}
\begin{lstlisting}[style=sample]
v1 v3 v2 v6 v4 v5
\end{lstlisting}
\end{taskbox}
\begin{mistakebox}
\begin{itemize}
  \item “编号小的顶点在前”不等于每次随便用普通队列；如果有多个入度为 0，要用最小堆。
  \item 边是有方向的，读入 a b 表示 a -> b，入度加在 b 上。
  \item 输出要带字母 v 前缀，不是只输出数字。
  \item 如果用 queue.Queue，无法保证同层编号最小优先。
  \item 邻接表下标要统一，顶点通常从 1 到 v。
\end{itemize}
\end{mistakebox}
\codesource
\begin{lstlisting}[style=py,caption={最终可用代码（对话来源）}]
# 这道题的重点是因为他需要在距离一致的时候用较小的序号插入，所以需要用堆排序

import heapq


def topoSort(G, n):
    # inDegree[i] 表示顶点 i 的入度
    # 这里顶点编号从 1 到 n，所以数组开 n + 1
    inDegree = [0] * (n + 1)

    # 统计每个顶点的入度
    for u in range(1, n + 1):
        for v in G[u]:
            inDegree[v] += 1

    # 小根堆，用来存当前所有入度为 0 的顶点
    # heapq 每次会弹出编号最小的顶点
    heap = []

    # 把一开始所有入度为 0 的顶点加入堆
    for i in range(1, n + 1):
        if inDegree[i] == 0:
            heapq.heappush(heap, i)

    # 保存拓扑排序结果
    seq = []

    # 只要堆不为空，就继续取点
    while heap:
        # 取出当前所有入度为 0 的点中编号最小的那个
        u = heapq.heappop(heap)

        # 加入拓扑序列
        seq.append(u)

        # 删除 u 指向的所有边
        for v in G[u]:
            inDegree[v] -= 1

            # 如果 v 的入度变成 0，说明它可以被选择了
            if inDegree[v] == 0:
                heapq.heappush(heap, v)

    # 如果没有输出所有顶点，说明图中有环
    if len(seq) != n:
        return None

    return seq


# 读入顶点数和弧数
n, a = map(int, input().split())

# G[u] 存放所有从 u 出发的点
# 顶点编号从 1 到 n
G = [[] for _ in range(n + 1)]

# 读入 a 条有向边
for _ in range(a):
    u, v = map(int, input().split())
    G[u].append(v)

# 求拓扑排序
ans = topoSort(G, n)

# 按题目要求输出 v1 v3 v2 ...
if ans is not None:
    print(" ".join("v" + str(x) for x in ans))
\end{lstlisting}
\begin{notebox}
\begin{itemize}
  \item 拓扑排序核心：入度数组 + 入度为 0 的集合。
  \item 要求编号小优先时，把队列换成最小堆。
  \item 输出格式常常比算法更容易扣分。
\end{itemize}
\end{notebox}


\subsection{2: 丛林中的路}
\begin{taskbox}
\textbf{题目编号与名称：}2: 丛林中的路\par
\textbf{要求：}给定若干村庄之间道路维护费用，要求保留一些道路使所有村庄连通且总费用最小，即求最小生成树总权值。\par
\textbf{输入：}多组数据；每组第一行 n，随后 n-1 行按字母给边；0 结束。\par
\textbf{输出：}每组输出最小维护费用。\par
\textbf{样例输入：}
\begin{lstlisting}[style=sample]
9
A 2 B 12 I 25
B 3 C 10 H 40 I 8
C 2 D 18 G 55
D 1 E 44
E 2 F 60 G 38
F 0
G 1 H 35
H 1 I 35
3
A 2 B 10 C 40
B 1 C 20
0
\end{lstlisting}
\textbf{样例输出：}
\begin{lstlisting}[style=sample]
216
30
\end{lstlisting}
\end{taskbox}
\begin{mistakebox}
\begin{itemize}
  \item 这是 MST，不是最短路；目标是让全图连通且总边权最小。
  \item 输入每行是变长格式：字母、k、再跟 k 组“字母 权值”。
  \item 村庄用 A,B,C... 编号，需要用 ord(ch)-ord("A") 转成下标。
  \item Kruskal 要对所有边按权值排序，并查集合并。
  \item 不要把无向边重复加两次影响不大，但本题输入本来只给后继村庄，按一条边处理即可。
\end{itemize}
\end{mistakebox}
\codesource
\begin{lstlisting}[style=py,caption={最终可用代码（对话来源）}]
import heapq
import sys

class Graph_adj_mat:

    def __init__(self, num_vertices):
        self.numVertices = num_vertices
        self.adj_matrix = [[0] * num_vertices for _ in range(num_vertices)]

    def addEdge(self, v1, v2, weight):
        self.adj_matrix[v1][v2] = weight
        self.adj_matrix[v2][v1] = weight   # 无向图

    def get_neighbors(self, idx):
        neighbors = []
        for j in range(self.numVertices):
            if self.adj_matrix[idx][j] != 0:
                neighbors.append(j)
        return neighbors

    def getWeight(self, v1, v2):
        return self.adj_matrix[v1][v2]


def prim(G, start):
    # pq 中存的是：边权、顶点编号
    pq = []

    # visited[i] 表示顶点 i 是否已经加入最小生成树
    visited = [False] * G.numVertices

    # 从起点 start 开始，代价为 0
    heapq.heappush(pq, (0, start))

    totalCost = 0

    while pq:
        # 取出当前边权最小的顶点
        currentCost, currentVert = heapq.heappop(pq)

        # 如果这个点已经加入生成树，就跳过
        if visited[currentVert]:
            continue

        # 把当前点加入最小生成树
        visited[currentVert] = True

        # 加上连接这个点的边权
        totalCost += currentCost

        # 遍历 currentVert 的所有邻居
        for nextVert in G.get_neighbors(currentVert):
            if not visited[nextVert]:
                newCost = G.getWeight(currentVert, nextVert)
                heapq.heappush(pq, (newCost, nextVert))

    return totalCost

while True:
    n = int(input())
    if n == 0:
        break

    G = Graph_adj_mat(n)

    for _ in range(n - 1):
        data = input().split()

        start = data[0]
        k = int(data[1])

        v1 = ord(start) - ord('A')

        index = 2
        for i in range(k):
            end = data[index]
            weight = int(data[index + 1])

            v2 = ord(end) - ord('A')

            G.addEdge(v1, v2, weight)

            index += 2
    m = prim(G,0)
    print(m)
\end{lstlisting}
\begin{notebox}
\begin{itemize}
  \item “连通且费用最小”基本就是最小生成树。
  \item Kruskal 模板：排序边 + 并查集。
  \item 字母编号题：A->0，B->1。
\end{itemize}
\end{notebox}


\subsection{3: A Walk Through the Forest}
\begin{taskbox}
\textbf{题目编号与名称：}3: A Walk Through the Forest\par
\textbf{要求：}办公室为 1，家为 2。无向带权图中，Jimmy 每一步必须“更接近家”：从 A 到 B 合法当且仅当 dist(B,2) < dist(A,2)。求从 1 到 2 的合法路线数。\par
\textbf{输入：}多组数据；每组第一行 N,M，之后 M 行 a b d；单独一行 0 结束。\par
\textbf{输出：}每组输出一行合法路线数。\par
\textbf{样例输入：}
\begin{lstlisting}[style=sample]
5 6
1 3 2
1 4 2
3 4 3
1 5 12
4 2 34
5 2 24
7 8
1 3 1
1 4 1
3 7 1
7 4 1
7 5 1
6 7 1
5 2 1
6 2 1
0
\end{lstlisting}
\textbf{样例输出：}
\begin{lstlisting}[style=sample]
2
4
\end{lstlisting}
\end{taskbox}
\begin{mistakebox}
\begin{itemize}
  \item 最短路要从家 2 出发跑 Dijkstra，得到每个点到家的最短距离。
  \item 路线计数不是最短路径条数，而是所有满足距离严格下降的路径数。
  \item 合法转移条件必须是 dist[v] < dist[u]，不能写成边权更小。
  \item DFS 需要 memo，否则会重复计算大量子问题。
  \item 多组输入第一行可能只有一个 0，要先判断长度或先读整行。
\end{itemize}
\end{mistakebox}
\codesource
\begin{lstlisting}[style=py,caption={最终可用代码（对话来源）}]
# 先建一个graph

import sys

import heapq

class Vertex:

    def __init__(self, key):
        self.id = key

        # connectedTo 用来保存当前顶点的所有邻居
        # 格式：
        # {
        #     邻居顶点对象1: 边权1,
        #     邻居顶点对象2: 边权2
        # }
        self.connectedTo = {}

        self.color = 'white'
        self.distance = float('inf')
        self.pred = None

    def addNeighbor(self, nbr, weight=0):
        # 添加一条从 self 到 nbr 的边，权重为 weight
        self.connectedTo[nbr] = weight

    def getConnections(self):
        # 返回所有邻居顶点对象
        return self.connectedTo.keys()

    def getId(self):
        return self.id

    def getWeight(self, nbr):
        # 返回 self 到 nbr 这条边的权重
        return self.connectedTo[nbr]

    def getPred(self):
        return self.pred

    def setPred(self, pred):
        self.pred = pred

    def getColor(self):
        return self.color

    def setColor(self, color):
        self.color = color

    def getDistance(self):
        return self.distance

    def setDistance(self, distance):
        self.distance = distance


class Graph:

    def __init__(self):
        # vertList 保存所有顶点
        # 格式：
        # {
        #     顶点编号: 顶点对象
        # }
        self.vertList = {}

        # 顶点数量
        self.numVertices = 0

    def addVertex(self, key):
        # 添加一个新顶点
        self.numVertices += 1
        newVertex = Vertex(key)
        self.vertList[key] = newVertex
        return newVertex

    def getVertex(self, key):
        # 根据编号获取顶点对象
        if key in self.vertList:
            return self.vertList[key]
        else:
            return None

    def __contains__(self, key):
        # 支持写法：if key in graph
        return key in self.vertList

    def addEdge(self, f, t, cost=0):
        # 如果起点 f 不存在，就先加入图中
        if f not in self.vertList:
            self.addVertex(f)

        # 如果终点 t 不存在，也先加入图中
        if t not in self.vertList:
            self.addVertex(t)

        # 添加一条从 f 到 t 的边
        self.vertList[f].addNeighbor(self.vertList[t], cost)

    def addUndirectedEdge(self, f, t, cost=0):
        # 无向图要加两条边
        # f -> t
        # t -> f
        self.addEdge(f, t, cost)
        self.addEdge(t, f, cost)

    def getVertices(self):
        # 返回所有顶点编号
        return self.vertList.keys()

    def __iter__(self):
        # 支持写法：for v in graph
        # 每次返回的是 Vertex 对象
        return iter(self.vertList.values())

def dijkstra(aGraph, start):
    # 初始化所有点
    for v in aGraph:
        v.setDistance(float('inf'))
        v.setPred(None)

    # 起点距离设为 0
    start.setDistance(0)

    # 小根堆中存放：
    # (当前距离, 顶点编号, 顶点对象)
    heap = []
    heapq.heappush(heap, (0, start.getId(), start))

    while heap:
        currentDist, _, currentVert = heapq.heappop(heap)

        # 如果这个距离不是最新的，跳过
        if currentDist != currentVert.getDistance():
            continue

        # 遍历当前顶点的所有邻居
        for nextVert in currentVert.getConnections():

            # 新路径长度 = 起点到 currentVert 的距离 + currentVert 到 nextVert 的边权
            newDist = currentVert.getDistance() + currentVert.getWeight(nextVert)

            # 如果新路径更短，就更新
            if newDist < nextVert.getDistance():
                nextVert.setDistance(newDist)
                nextVert.setPred(currentVert)

                heapq.heappush(heap, (newDist, nextVert.getId(), nextVert))

def dfs_count_routes(u, memo):
    if u.getId() == 2:
        return 1

    if memo[u.getId()] != -1:
        return memo[u.getId()]

    ans = 0

    for v in u.getConnections():
        if v.getDistance() < u.getDistance():
            ans += dfs_count_routes(v, memo)

    memo[u.getId()] = ans
    return ans

sys.setrecursionlimit(1000000)

while True:
    line = sys.stdin.readline().strip()

    if line == "0":
        break

    n, m = map(int, line.split())

    G = Graph()

    # 先创建 1 到 n 的所有顶点
    for i in range(1, n + 1):
        G.addVertex(i)

    for _ in range(m):
        a, b, d = map(int, sys.stdin.readline().split())
        G.addUndirectedEdge(a, b, d)
    
    dijkstra(G, G.getVertex(2))

    memo = [-1] * (n + 1)

    print(dfs_count_routes(G.getVertex(1), memo))
    

\end{lstlisting}
\begin{notebox}
\begin{itemize}
  \item 先算“到家距离”，再沿距离下降方向 DP。
  \item return memo[u] 表示之前算过；return ans 表示这次刚算完。
  \item 这题合法路线数量，不是最短路线数量。
\end{itemize}
\end{notebox}


\subsection{4: 打怪救公主}
\begin{taskbox}
\textbf{题目编号与名称：}4: 打怪救公主\par
\textbf{要求：}迷宫中从 * 出发救 +，数字格是怪物血量，进入会扣血；\# 是墙。若能到达公主，输出剩余最大血量，否则 0。\par
\textbf{输入：}第一行 m,n,t；接下来 m 行地图。\par
\textbf{输出：}能救出则输出到达 + 时最大剩余血量，否则输出 0。\par
\textbf{样例输入：}
\begin{lstlisting}[style=sample]
5 6 10
..*...
.#2###
5#..4#
.##9.#
.#+..#
\end{lstlisting}
\textbf{样例输出：}
\begin{lstlisting}[style=sample]
4
\end{lstlisting}
\end{taskbox}
\begin{mistakebox}
\begin{itemize}
  \item 这题不是普通 BFS 求最少步数，而是要最大化剩余血量。
  \item 同一个格子如果以更多血量到达，仍然值得继续扩展；visited 不能只是布尔值。
  \item 可以用 best[x][y] 记录到达该格子的最大剩余血量，只有更优才入队。
  \item 进入数字格时先扣血，血量必须仍然大于 0 才能继续。
  \item 起点 * 和终点 + 不扣血，. 也不扣血。
\end{itemize}
\end{mistakebox}
\codesource
\begin{lstlisting}[style=py,caption={最终可用代码（对话来源）}]
from collections import deque

def max_blood_bfs(maze, M, N, sx, sy, B):
    q = deque()
    max_blood = [[-1]*N for _ in range(M)]
    q.append((sx, sy, B))
    max_blood[sx][sy] = B

    directions = [(-1,0),(1,0),(0,-1),(0,1)]

    while q:
        x, y, blood = q.popleft()

        for dx, dy in directions:
            nx, ny = x+dx, y+dy
            if nx<0 or nx>=M or ny<0 or ny>=N:
                continue
            cell = maze[nx][ny]
            if cell == '#':
                continue

            new_blood = blood
            if cell.isdigit():
                dmg = int(cell)
                new_blood -= dmg
                if new_blood <=0:
                    continue

            if new_blood > max_blood[nx][ny]:
                max_blood[nx][ny] = new_blood
                q.append((nx, ny, new_blood))

    # 找到公主位置
    for i in range(M):
        for j in range(N):
            if maze[i][j] == '+':
                return max_blood[i][j] if max_blood[i][j] != -1 else 0

M, N, B = map(int, input().split())
maze = []
sx, sy = -1, -1

for i in range(M):
    row = input().strip()
    maze.append(row)

    for j in range(N):
        if row[j] == '*':
            sx, sy = i, j

print(max_blood_bfs(maze, M, N, sx, sy,B))
\end{lstlisting}
\begin{notebox}
\begin{itemize}
  \item 看到“剩余最大血量”，visited 要存最优值，不是 True/False。
  \item 只有状态变好才继续入队。
  \item 血量减到 0 会死，因此必须 nhp > 0。
\end{itemize}
\end{notebox}


\subsection{5: 兔子与樱花}
\begin{taskbox}
\textbf{题目编号与名称：}5: 兔子与樱花\par
\textbf{要求：}给定地点、无向带权道路和若干查询，输出每个查询的最短路径，并按 A->(d)->B 的格式打印整条路线。\par
\textbf{输入：}先读 P 个地点；再读 Q 条道路；最后读 R 个查询。\par
\textbf{输出：}每个查询输出一行最短路径表示；起点等于终点时只输出地点名。\par
\textbf{样例输入：}
\begin{lstlisting}[style=sample]
6
Ginza
Sensouji
Shinjukugyoen
Uenokouen
Yoyogikouen
Meijishinguu
6
Ginza Sensouji 80
Shinjukugyoen Sensouji 40
Ginza Uenokouen 35
Uenokouen Shinjukugyoen 85
Sensouji Meijishinguu 60
Meijishinguu Yoyogikouen 35
2
Uenokouen Yoyogikouen
Meijishinguu Meijishinguu
\end{lstlisting}
\textbf{样例输出：}
\begin{lstlisting}[style=sample]
Uenokouen->(35)->Ginza->(80)->Sensouji->(60)->Meijishinguu->(35)->Yoyogikouen
Meijishinguu
\end{lstlisting}
\end{taskbox}
\begin{mistakebox}
\begin{itemize}
  \item 不仅要求最短距离，还要求输出具体路径，因此需要记录前驱。
  \item 地点是字符串，先映射成编号，输出时再映射回名称。
  \item 多次查询且 P<30，可以用 Floyd 预处理所有点对最短路。
  \item 路径恢复时要维护 next 矩阵，而不是只保存 dist。
  \item 起点等于终点时，直接输出地点名，不要输出箭头。
\end{itemize}
\end{mistakebox}
\codesource
\begin{lstlisting}[style=py,caption={最终可用代码（对话来源）}]
class Vertex:
    def __init__(self, key, value):
        self.id = key
        self.name = value
        self.connectedTo = {}
        self.distance = float('inf')
        self.pred = None
 
    def addNeighbor(self, nbr, weight=0):
        self.connectedTo[nbr] = weight

    def getConnections(self):
        return self.connectedTo.keys()

    def getId(self):
        return self.name

    def getWeight(self, nbr):
        return self.connectedTo[nbr]

    def getDistance(self):
        return self.distance

    def setDistance(self, distance):
        self.distance = distance

    def getPred(self):
        return self.pred

    def setPred(self, pred):
        self.pred = pred


class Graph:
    def __init__(self):
        self.vertList = {}
        self.nameMap = {}
        self.numVertices = 0

    def addVertex(self, key, name):
        newVertex = Vertex(key, name)
        self.vertList[key] = newVertex
        self.nameMap[name] = newVertex
        self.numVertices += 1
        return newVertex

    def getVertexByName(self, name):
        return self.nameMap.get(name, None)

    def addEdgeByName(self, f_name, t_name, cost=0):
        f = self.getVertexByName(f_name)
        t = self.getVertexByName(t_name)
        f.addNeighbor(t, cost)

    def __iter__(self):
        return iter(self.vertList.values())

    def reset(self):
        for v in self:
            v.setDistance(float('inf'))
            v.setPred(None)

def dijkstra_no_pq(aGraph, start):
    aGraph.reset()
    # 设置起点到自己的距离为 0
    start.setDistance(0)

    # 已经确定最短路径的顶点集合
    visited = set()

    # 只要还有未访问顶点，就继续循环
    while len(visited) < aGraph.numVertices:

        # 在所有未访问顶点中，找到距离起点最小的顶点
        currentVert = None
        currentDist = float('inf')
        for v in aGraph:
            if v not in visited and v.getDistance() < currentDist:
                currentDist = v.getDistance()
                currentVert = v

        # 如果没有可访问顶点，说明剩余顶点不可达，结束循环
        if currentVert is None:
            break

        # 把当前顶点加入已访问集合
        visited.add(currentVert)

        # 遍历 currentVert 的所有邻接点
        for nextVert in currentVert.getConnections():

            # 计算通过 currentVert 到 nextVert 的新距离
            newDist = currentVert.getDistance() + currentVert.getWeight(nextVert)

            # 如果新距离更短，更新 nextVert 的最短距离
            if newDist < nextVert.getDistance():
                nextVert.setDistance(newDist)
                nextVert.setPred(currentVert)
def getPath(end):
    path = []

    current = end
    while current is not None:
        path.append(current)
        current = current.getPred()

    path.reverse()
    return path

def printPath(path):
    ans = path[0].getId()

    for i in range(1, len(path)):
        weight = path[i-1].getWeight(path[i])
        ans += "->(" + str(weight) + ")->" + path[i].getId()

    print(ans)

g = Graph()
n1 = int(input())

for i in range(n1):
    name = input().strip()
    g.addVertex(i, name)

n2 = int(input())
for _ in range(n2):
    f_name, t_name, weight = input().split()
    weight = int(weight)
    g.addEdgeByName(f_name, t_name, weight)
    g.addEdgeByName(t_name, f_name, weight)

n3 = int(input())
for _ in range(n3):
    s_name, e_name = input().split()
    s = g.getVertexByName(s_name)
    e = g.getVertexByName(e_name)
    dijkstra_no_pq(g, s)
    p = getPath(e)
    printPath(p)
\end{lstlisting}
\begin{notebox}
\begin{itemize}
  \item 多源多查询且点数小：Floyd 很方便。
  \item 输出路径需要 next 矩阵恢复下一跳。
  \item 格式是 点名->(边长)->点名，不是只输出总距离。
\end{itemize}
\end{notebox}


\subsection{56: 拯救行动}
\begin{taskbox}
\textbf{题目编号与名称：}56: 拯救行动\par
\textbf{要求：}在 $N\times M$ 牢房矩阵中，骑士 r 要到达公主 a。道路为 @，墙为 \#，守卫为 x。骑士上下左右移动，每移动一格花费 1；如果进入守卫格，还需要额外花费 1 杀死守卫。求成功救到公主的最短时间；不可达输出 Impossible。\par
\textbf{输入：}第一行整数 S，表示数据组数。每组先输入 N 和 M，随后 N 行、每行 M 个字符，字符包括 @、\#、x、r、a。\par
\textbf{输出：}每组输出最短时间；若不可能成功，输出 Impossible。\par
\textbf{样例输入：}
\begin{lstlisting}[style=sample]
2
7 8
#@#####@
#@a#@@r@
#@@#x@@@
@@#@@#@#
#@@@##@@
@#@@@@@@
@@@@@@@@ 
13 40
@x@@##x@#x@x#xxxx##@#x@x@@#x#@#x#@@x@#@x
xx###x@x#@@##xx@@@#@x@@#x@xxx@@#x#x@@x@
#@x#@x#x#@@##@@x#@xx#xxx@@x##@@@#@x@@x@x
@##x@@@x#xx#@@#xxxx#@@x@x@#@x@@@x@#@#x@#
@#xxxxx##@@x##x@xxx@@#x@x####@@@x#x##@#@
#xxx#@#x##xxxx@@#xx@@@x@xxx#@#xxx@x#####
#x@xxxx#@x@@@@##@x#xx#xxx@#xx#@#####x#@x
xx##@#@x##x##x#@x#@a#xx@##@#@##xx@#@@x@x
x#x#@x@#x#@##@xrx@x#xxxx@##x##xx#@#x@xx@
#x@@#@###x##x@x#@@#@@x@x@@xx@@@@##@@x@@x
x#xx@x###@xxx#@#x#@@###@#@##@x#@x@#@@#@@
#@#x@x#x#x###@x@@xxx####x@x##@x####xx#@x
#x#@x#x######@@#x@#xxxx#xx@@@#xx#x#####@
\end{lstlisting}
\textbf{样例输出：}
\begin{lstlisting}[style=sample]
13
7
\end{lstlisting}
\end{taskbox}
\begin{mistakebox}
\begin{itemize}
  \item 这题不是普通 BFS，因为进入不同格子的代价不一样：普通道路代价 1，守卫 x 代价 2。所以如果直接用队列按层扩展，可能先到达的状态并不是最短时间。
  \item 正确思路应当是 Dijkstra：优先队列里存当前最短时间，每次弹出时间最小的位置。
  \item 不能把 Step 类对象直接放入 heapq 比较；如果堆元素是对象，需要定义 \texttt{\_\_lt\_\_}，更简单的写法是直接压入 \texttt{(time, x, y)} 元组。
  \item 距离数组 \texttt{dist[x][y]} 要记录到每个格子的最短时间。只有当新时间更短时，才更新并入堆。
  \item 读入地图时要顺便记录 r 和 a 的位置；多组数据时，每组都要重新初始化 maze、dist、pq。
  \item 遇到 \# 必须跳过；@、r、a 的移动代价都是 1，x 的移动代价是 2。
  \item 可以在 Dijkstra 弹出公主位置时直接返回，因为这时已经保证是全局最短时间。
\end{itemize}
\end{mistakebox}
\codesource
\begin{lstlisting}[style=py,caption={最终可用代码（对话来源）}]
import heapq


n = int(input())

for _ in range(n):

    maze = []

    M,N = map(int,input().split())

    # 找起点和终点
    for i in range(M):
        row = input().strip()
        maze.append(row)
        for j in range(N):
            if maze[i][j] == 'r':
                sx, sy = i, j
            elif maze[i][j] == 'a':
                ex, ey = i, j

    # 优先队列 (steps, x, y)
    pq = []
    heapq.heappush(pq, (0, sx, sy))

    # 记录每个格子最短时间
    dist = [[float('inf')] * N for _ in range(M)]
    dist[sx][sy] = 0

    directions = [(-1,0), (1,0), (0,-1), (0,1)]

    found = False

    while pq:
        steps, x, y = heapq.heappop(pq)

        # 如果当前步数比记录的还大，说明已经有更短路径了
        if steps > dist[x][y]:
            continue

        # 到达终点
        if x == ex and y == ey:
            print(steps)
            found = True
            break

        # 扩展邻居
        for dx, dy in directions:
            nx = x + dx
            ny = y + dy

            if 0 <= nx < M and 0 <= ny < N:
                if maze[nx][ny] == '#':
                    continue

                # 走普通格子: +1, 守卫: +2
                if maze[nx][ny] == 'x':
                    nt = steps + 2
                else:
                    nt = steps + 1

                # 如果这条路径更短，就加入队列
                if nt < dist[nx][ny]:
                    dist[nx][ny] = nt
                    heapq.heappush(pq, (nt, nx, ny))

    if not found:
        print("Impossible")
\end{lstlisting}
\begin{notebox}
\begin{itemize}
  \item 网格题如果每一步代价完全相同，用 BFS；如果边权出现 1 和 2 这种差异，用 Dijkstra。
  \item \texttt{heapq} 最稳的写法是堆里放元组：\texttt{(距离, 行, 列)}。
  \item Dijkstra 的过期状态判断：弹出后如果 \texttt{t != dist[x][y]}，说明它已经不是最新最优状态，直接跳过。
\end{itemize}
\end{notebox}


\subsection{57: 连接所有点的最小费用}
\begin{taskbox}
\textbf{题目编号与名称：}57: 连接所有点的最小费用\par
\textbf{要求：}给定二维平面上的若干点，连接任意两点的费用为曼哈顿距离 $|x_i-x_j|+|y_i-y_j|$。要求用最小总费用把所有点连通，并且任意两点之间有且仅有一条简单路径，即求完全图上的最小生成树总权重。\par
\textbf{输入：}一行，一个 points 数组，例如 \texttt{[[0,0],[2,2],[3,10],[5,2],[7,0]]}。可用 \texttt{ast.literal\_eval} 读成二维列表。\par
\textbf{输出：}一个整数，表示连接所有点的最小总费用。\par
\textbf{样例输入：}
\begin{lstlisting}[style=sample]
[[0,0],[2,2],[3,10],[5,2],[7,0]]
\end{lstlisting}
\textbf{样例输出：}
\begin{lstlisting}[style=sample]
20
\end{lstlisting}
\end{taskbox}
\begin{mistakebox}
\begin{itemize}
  \item 题目说“所有点连通，且任意两点之间有且仅有一条简单路径”，本质就是最小生成树 MST，不是最短路。
  \item 每两个点之间都可以连边，所以这是一个完全图；没有必要真的把所有边都存成邻接表再排序，直接用 Prim 更省空间。
  \item 曼哈顿距离要写成 \texttt{abs(x1-x2)+abs(y1-y2)}，不要误用欧氏距离。
  \item 输入不是普通的数字矩阵，而是一整行 Python 风格列表；要用 \texttt{ast.literal\_eval}，不能直接 \texttt{split()}。
  \item Prim 中 \texttt{minDist[i]} 表示“当前已选集合到点 i 的最小连接费用”，每次选未加入集合且 \texttt{minDist} 最小的点。
  \item 每加入一个新点后，要用它去更新所有未加入点的 \texttt{minDist}。
  \item 若点数较小，$O(n^2)$ Prim 比建 $O(n^2)$ 条边再 Kruskal 更直接，也避免内存浪费。
\end{itemize}
\end{mistakebox}
\codesource
\begin{lstlisting}[style=py,caption={最终可用代码（对话来源）}]
import ast

import heapq
import sys


class Vertex:
    def __init__(self, key):
        self.id = key
        self.connectedTo = {}
        self.color = 'white' # 顶点访问状态（white=未访问，gray=正在访问，black=已完成）
        self.distance = float('inf') # 当前顶点到起始点的距离，初始设为无穷大
        self.pred = None # 处理这个顶点的前驱，用于记录路径
 
    def addNeighbor(self, nbr, weight=0):
        self.connectedTo[nbr] = weight # 是字典，nbr是顶点对象的key

    def __str__(self):
        return str(self.id) + ' connectedTo: ' + str([x.id for x in self.connectedTo])

    def getConnections(self):
        return self.connectedTo.keys()

    def getId(self):
        return self.id

    def getWeight(self, nbr):
        return self.connectedTo[nbr]
    
    def setDistance(self, dist):
        self.distance = dist

    def getDistance(self):
        return self.distance

    def setPred(self, pred):
        self.pred = pred

class Graph:

    def __init__(self):
        self.vertList = {}
        self.numVertices = 0

    def addVertex(self, key):
        self.numVertices = self.numVertices + 1
        newVertex = Vertex(key)
        self.vertList[key] = newVertex
        return newVertex

    def getVertex(self, n): # 通过key查找顶点
        if n in self.vertList:
            return self.vertList[n]
        else:
            return None

    def __contains__(self, n):
        return n in self.vertList
    
    def addEdge(self, f, t, cost=0): # f起点，t终点
        if f not in self.vertList:
            nv = self.addVertex(f)

        if t not in self.vertList:
            nv = self.addVertex(t)
        self.vertList[f].addNeighbor(self.vertList[t], cost)

    def getVertices(self):
        return self.vertList.keys()

    def __iter__(self):
        return iter(self.vertList.values())


def prim(G, start):
    # heapq 实现的优先队列
    pq = []

    # 记录哪些顶点已经加入最小生成树
    inTree = set()

    # 初始化所有顶点
    for v in G:
        v.setDistance(sys.maxsize)
        v.setPred(None)

    # 起点距离设为 0
    start.setDistance(0)

    # 修改：heapq 里面放三元组
    # (当前最小边权, 顶点编号, 顶点对象)
    # 加顶点编号是为了防止两个 Vertex 对象不能比较而 Runtime Error
    heapq.heappush(pq, (0, start.getId(), start))

    totalCost = 0

    while pq:
        # 相当于 pq.delMin()
        currentDistance, currentId, currentVert = heapq.heappop(pq)

        # 修改：如果这个点已经加入生成树，就跳过
        # 因为 heapq 没有 decreaseKey，旧的较大距离可能还留在堆里
        if currentVert in inTree:
            continue

        # 当前顶点正式加入最小生成树
        inTree.add(currentVert)

        # 当前边权加入答案
        totalCost += currentDistance

        # 遍历邻接点
        for nextVert in currentVert.getConnections():

            newCost = currentVert.getWeight(nextVert)

            # 修改：原来是 if nextVert in pq
            # heapq 不能直接判断某个点是否还在队列中
            # 所以改成判断 nextVert 是否还没加入生成树
            if nextVert not in inTree and newCost < nextVert.getDistance():

                nextVert.setPred(currentVert)

                nextVert.setDistance(newCost)

                # 修改：heapq 没有 decreaseKey
                # 所以直接把新的更小距离重新压入堆
                heapq.heappush(pq, (newCost, nextVert.getId(), nextVert))

    return totalCost



input_str = input().strip()
points = ast.literal_eval(input_str)

g = Graph()

for i in range(len(points)):
    for j in range(i, len(points)):
        w = abs(points[i][0] - points[j][0]) + abs(points[i][1] - points[j][1])
        g.addEdge(i,j,w)
        g.addEdge(j,i,w)

result = prim(g, g.getVertex(0))

print(result)
\end{lstlisting}
\begin{notebox}
\begin{itemize}
  \item 看到“把所有点/村庄/节点连接起来，总代价最小”就优先想到 MST。
  \item 点之间天然都有边、边权可即时计算时，Prim 常比 Kruskal 更方便。
  \item 这题和“丛林中的路”同属最小生成树：前者是坐标完全图，后者是显式给边图。
\end{itemize}
\end{notebox}



\section{全局输入与调试清单}
\begin{longtable}{p{0.23\textwidth}p{0.69\textwidth}}
\toprule
\textbf{检查项} & \textbf{考前复盘说明} \\
\midrule
读题先看输出口径 & 最短步数、经过格子数、序列长度、剩余血量、路线数量是不同目标。仙岛求药计数包括起点，所以从 1 开始；变换迷宫输出时间，所以从 0 开始。\\
多组数据结束条件 & 常见结束方式有 EOF、0、0 0、先给 T。A Walk Through the Forest 是单独一行 0 结束；二叉树编号子树是 0 0 结束；前中序建树是 EOF。\\
行列顺序 & 网格题统一使用 maze[row][col]，也就是 x 表示行、y 表示列。visited 的大小必须是 [行数][列数]，不要写成 [N][N]。\\
visited 是布尔还是最优值 & 普通最短路 BFS 用 bool visited；动态迷宫要 visited[x][y][t\%K]；打怪救公主要 best[x][y] 存最大剩余血量。\\
何时标记 visited & BFS 通常在入队时标记，避免同一个状态被重复加入队列。若状态可能被更优值更新，则不能用布尔 visited。\\
边界判断顺序 & 先判断越界，再判断墙/障碍，再判断是否访问过。否则容易数组越界 RE。\\
起点距离初始化 & 题目说“经过的方格数包括初始位置”，起点距离为 1；题目说“花费时间”，起点时间为 0。\\
图是有向还是无向 & 拓扑排序是有向图；丛林中的路、A Walk Through the Forest、兔子与樱花都是无向图。无向图读边时两边都要加。\\
最短路还是最小生成树 & “从 A 到 B 最短”是最短路；“保证所有点连通且总费用最小”是最小生成树。不要把 Prim/Kruskal 和 Dijkstra 混用。\\
编号从 0 还是 1 开始 & 题目顶点编号若为 1..n，就开 n+1 数组；字母编号 A..Z 用 ord(ch)-ord('A')。\\
输出格式 & 是否有前缀 v、是否有箭头、是否有空格、是否末尾无空格，都要按样例。拓扑排序输出 v1 v3；兔子与樱花输出 A->(d)->B。\\
递归出口 & 树题的空结点通常返回空串或 0；扩展二叉树遇到 . 返回 None；编号为 -1 的孩子深度为 0。\\
树的高度口径 & 按结点数：空树 0、叶子 1。按边数：根高度 0。题目 49 按结点数；题目 27 转换高度按样例可知是边数口径。\\
字符串建树切分 & 前序+中序：根是 pre[0]；中序+后序：根是 post[-1]。切分另一个序列时一定用左子树长度。\\
heapq 常见拼写 & import heapq；heapq.heapify(h)；heapq.heappush(h,x)；heapq.heappop(h)。不要写 headq 或 heapop。\\
Dijkstra 过期节点 & 弹出 (d,u) 后若 d != dist[u]，说明是旧状态，要 continue。\\
DFS + 记忆化 & A Walk Through the Forest 的路线计数必须 memo。return memo[u] 是复用旧答案；return ans 是本次刚算完的结果。\\
\bottomrule
\end{longtable}

\section{考前模板索引}
\begin{longtable}{p{0.24\textwidth}p{0.22\textwidth}p{0.46\textwidth}}
\toprule
\textbf{题面关键词} & \textbf{模板} & \textbf{关键判断} \\
\midrule
最少步数、最少方格数 & 普通 BFS & 队列里放 (x,y,dist)，visited 二维；若包括起点，dist 从 1 开始。\\
迷宫随时间变化 & 状态 BFS & 状态必须包含时间模数，例如 (x,y,t\%K)。\\
同一位置可用更好资源重访 & 最优值 BFS/SPFA 思想 & visited 改成 best 数组，只有更优状态才入队。\\
每次只能改变一个字母 & 单词图 BFS & 单词作为点，只差一个字母连边，起点距离为序列长度 1。\\
同等条件编号小优先 & 堆优化拓扑 & 入度为 0 的点放最小堆，不用普通队列。\\
所有点连通且费用最小 & MST & Kruskal：边排序 + 并查集；Prim：点扩展 + 堆。\\
从某点到所有点最短 & Dijkstra & 非负权图；dist 初始化 INF，起点为 0，堆中弹最小。\\
多源多查询最短路 & Floyd & 点数小，用三重循环，并用 next 矩阵恢复路径。\\
合法路线数量 & 最短路 + DP & 先算评价函数 dist，再按 dist 严格下降方向 DFS 记忆化。\\
前序 + 中序建树 & 字符串递归 & root=pre[0]；中序划分；输出后序 left+right+root。\\
中序 + 后序建树 & 字符串递归 & root=post[-1]；中序划分；输出前序 root+left+right。\\
扩展二叉树先序 & 递归读入 & 遇到 . 返回空；先建左再建右。\\
括号嵌套树 & 括号深度解析 & 根是第一个字符；只按最外层逗号切左右子树。\\
BST 插入和层序 & BST + deque & 重复值忽略；层序用 deque.popleft。\\
验证 BST & 范围递归 & 每个结点携带 (low, high)，不能只看父子大小。\\
Huffman/WPL & 最小堆 & 每次取两个最小值合并，答案加合并和。\\
完全二叉树编号 & 区间扩展 & 子树每一层是连续编号区间 [left,right]。\\
\bottomrule
\end{longtable}

\subsection*{最小可背 BFS 模板}
\begin{lstlisting}[style=py]
from collections import deque
q = deque([(sx, sy, 0)])
visited = [[False] * C for _ in range(R)]
visited[sx][sy] = True
while q:
    x, y, d = q.popleft()
    for dx, dy in [(-1,0),(1,0),(0,-1),(0,1)]:
        nx, ny = x + dx, y + dy
        if nx < 0 or nx >= R or ny < 0 or ny >= C:
            continue
        if visited[nx][ny] or grid[nx][ny] == '#':
            continue
        visited[nx][ny] = True
        q.append((nx, ny, d + 1))
\end{lstlisting}

\subsection*{最小可背 Dijkstra 模板}
\begin{lstlisting}[style=py]
import heapq
INF = 10**18
dist = [INF] * (n + 1)
dist[start] = 0
pq = [(0, start)]
while pq:
    d, u = heapq.heappop(pq)
    if d != dist[u]:
        continue
    for v, w in graph[u]:
        nd = d + w
        if nd < dist[v]:
            dist[v] = nd
            heapq.heappush(pq, (nd, v))
\end{lstlisting}

\subsection*{最小可背并查集模板}
\begin{lstlisting}[style=py]
parent = list(range(n))

def find(x):
    if parent[x] != x:
        parent[x] = find(parent[x])
    return parent[x]

def union(a, b):
    ra, rb = find(a), find(b)
    if ra == rb:
        return False
    parent[rb] = ra
    return True
\end{lstlisting}

\end{document}
